\documentclass[11pt,a4paper]{article}
\usepackage[margin=1.15in]{geometry}
\usepackage[T1]{fontenc}
\usepackage{amsmath,amssymb}
\usepackage{booktabs}
\usepackage[expansion=false]{microtype}
\usepackage[hidelinks]{hyperref}
\usepackage{textcomp}

\newcommand{\game}{\textsc{The League}}
\title{\textsc{The League}:\\[2pt] \large Technical Notes on a Game of University Competition under Two Institutions\thanks{The game, the simulations, and a first draft of these notes were produced in collaboration with Claude (Anthropic). The game is a single, dependency-free HTML file; it, the simulation harness, and the source of these notes live at \url{https://github.com/steffenhuck/ivy}. The design owes several of its instruments, and one permanent benchmark strategy, to a playtest by Rustam Hakimov. Comments are welcome, and so are high scores.}}
\author{Steffen Huck\\ \normalsize WZB Berlin \& University College London}
\date{September 2026}

\begin{document}
\maketitle

\begin{abstract}
\noindent Four universities compete for forty school-leavers, once a year, for twenty years. Seats cost rent in both worlds, filled or not; in one world colleges set their own fees and must honour every acceptance, in the other a regulator fixes the fee and a clearing house assigns students by deferred acceptance against reported quotas. \game{} is a browser game in which a player takes over the worst of the four under either institution and tries to reach the top of a published league table. These notes set out the common environment, the two market designs, and a calibration by scripted players of three skill levels, with a full-spending benchmark and probe strategies patrolling the design's edges. The market kills by crowd; the scheme kills by emptiness.
\end{abstract}

\section{Introduction}

You have accepted the worst chair in academia. Greyfriars College is fourth of four in the league table that \emph{The Morning Ledger} publishes every autumn; its endowment would not survive two bad roofs; its teaching is tired and its research is a rumour. The governing board, out of options and nearly out of money, asks only that you climb.

In either world each of your two departments keeps eight seats, and every seat costs rent whether a student sits in it or not. How you may climb depends on what else the world permits. In the first world you set your own fees. Applicants apply wherever they can afford, you make an offer to every applicant above a bar of your choosing, and every offer that is accepted must be honoured---the ninth student in a department, the one without a seat, costs a small fortune in emergency provision. In the second world the fee has been confiscated by a regulator and set equal for everyone. A national clearing house runs the admissions round: you report a quota and a standard, an algorithm assigns each student to the best place that will hold them, and the quota is never exceeded. The market can overflow your seats; the scheme can only leave them empty.

\game{} is a browser game in which both worlds can be played, from the same starting position, against the same opponents. These notes are its documentation, written for economists.

The design brief imposed a discipline worth recording. The engine was built and calibrated headless before any interface existed: three scripted players of increasing sophistication---static, adaptive, exploitative---were run over five hundred seeded histories per institution, and parameters were tuned until the sophisticated player usually carries the worst college to the top, the adaptive player improves but rarely wins, and the static player does not improve at all; a full-spending benchmark and a set of probe strategies---an overbooker, a scholarship spammer, a card gambler---patrol the result (Section~\ref{sec:calib}). The calibration tables are reproduced in Section~\ref{sec:calib} and shipped inside the game file itself. A second discipline is pedagogical. In each world the winning strategy, discoverable by play, is a known piece of economics: in the market it is monopoly pricing over a captive segment whose demand curve must be learned by costly experiment; in the scheme it is the arithmetic of a stable mechanism---what a quota report does and does not buy (S\"onmez, 1997), and what it is worth against rivals who spend rationally. The syllabus is smuggled in with the fun.

The notes proceed as follows. Section~\ref{sec:env} sets out the environment common to both worlds. Sections~\ref{sec:market} and~\ref{sec:scheme} specify the two institutions. Section~\ref{sec:calib} reports the calibration and a comparison of the two games' difficulty. Section~\ref{sec:disc} discusses what the comparison teaches.

\section{The common environment}\label{sec:env}

\paragraph{Universities.} There are four universities, indexed $i$. Each carries five state variables: an endowment $E^i$ and four qualities---research and teaching in each of two fields, STEM ($S$) and the humanities and social sciences ($H$)---written $R_S^i, T_S^i, R_H^i, T_H^i$. All money is in thousands of pounds. The league table ranks universities by the sum of the four qualities; ties are broken by endowment, then by a fixed index. The table is public. Endowments, fees, thresholds, quotas, and enrolments are not. One university is run by the player; the other three are run by heuristic personalities described in Appendix~\ref{app:ai}. Starting states appear in Table~\ref{tab:start}.

\paragraph{Students.} Each year forty fresh applicants are drawn. An applicant is a quadruple $(s, b, \theta, f)$. The school score is $s \sim N(60, 13^2)$, truncated to $[20,100]$. The budget is $b = 10 + 0.12\,(s - 60) + \varepsilon$ with $\varepsilon \sim N(0, 3.8^2)$, truncated to $[3,25]$; money and marks travel together, mildly, and the truncation point equals the fee cap, so nothing but demand polices a high fee. The taste for research is $\theta \sim U[0, 1.2]$. The field is assigned by stratification: exactly $\mathrm{round}(40\,p_t)$ applicants prefer STEM, where $p_1 = \tfrac12$ and
\[
p_{t+1} \;=\; \mathrm{clamp}\!\left(p_t + 0.02\,\tanh\!\Big(\tfrac{\overline{R}_{S,t} - \overline{R}_{H,t}}{10}\Big)\right)\!,
\qquad p_t \in [0.3, 0.7],
\]
so that fashion drifts gently toward whichever field the sector researches better. A student attending university $i$ in field $f$ enjoys utility $u^i = T_f^i + \theta R_f^i$: teaching benefits everyone alike, research prestige benefits those who care for it. Students live one year. They pay one fee, are taught once, and leave; there are no cohorts.

\paragraph{Quality dynamics.} Investment has diminishing returns and takes effect the following year. For each quality $Q \in \{R_S, T_S, R_H, T_H\}$ with investment $I_Q \ge 0$,
\[
R' \;=\; \delta R + \gamma \sqrt{I_R}, \qquad
T' \;=\; \delta T + \gamma \sqrt{I_T} + \kappa\,(\bar{s} - 60)\cdot\mathbf{1}\{\text{anyone admitted}\},
\]
with decay $\delta = 0.87$, productivity $\gamma = 1.5$, and intake effect $\kappa = 0.14$, where $\bar{s}$ is the department's mean admitted score. Qualities are floored at zero. Decay is deliberately brutal: an unattended quality loses thirteen per cent a year, so standing still is a policy with consequences.

\paragraph{The cost of seats.} Both worlds share one cost function. A department with enrolment $m$ pays
\[
C \;=\; 8\rho \;+\; c\,\max(0,\, m - 8), \qquad \rho = 0.75, \quad c = 24:
\]
rent on all eight seats, filled or not, plus a steep price per student beyond them. Deferred acceptance never over-fills a quota, so the second term is reachable only in the market. Fees may exceed $c$ (the cap is 25): overflow-for-profit is priced, not forbidden---and, as Section~\ref{sec:calib} reports, never observed.

\paragraph{The ledger.} Fees are paid up front. With fee income $F_t$, the year's bills $C_t$---seat costs as above plus any stipend bill---and investment $I_t$,
\[
E_{t+1} \;=\; \big(E_t + F_t - C_t - I_t\big)\,(1 + r), \qquad r = 0.05 .
\]
Bankruptcy is checked after admissions and before spending: if $E_t + F_t - C_t < 0$, the college is finished. For the player this ends the game. A broke rival is frozen: it stops admitting, invests nothing, and decays.

\paragraph{Merit scholarships.} In either world a college sets, per field and per year, a stipend $\sigma \in [0, 20]$ per student and a qualifying bar $\beta \in [40, 95]$ in school score. Every enrolling student with $s \ge \beta$ is owed the stipend, so the bill is $\sigma$ times the qualifying intake---known only after clearing, part of $C_t$, and therefore able to bankrupt: a generous stipend at a low bar is a road to ruin, not a bargain, and the calibration probes confirm it (Section~\ref{sec:calib}). The instrument is deliberately per-head rather than a fund: a student asked to value a fund must reason about its division among rivals she cannot observe, while a student offered a stipend needs to know only her own offer.

The stipend's channel differs by world, on principle: cash nets against a price where a price exists. In the market it counts toward the fee---a qualifying student can afford the college if $\varphi_f - \sigma \le b$---while choice among offers stays quality-only, exactly as fees themselves never sway choice in this model; merit aid is thus a targeted price cut to the bright, and its work is done at the gate. In the scheme there is no price to net against, so the stipend is a side payment entering preferences: a qualifying student ranks the paying college $\sigma$ utility points higher (one point per thousand), cash being cash even where fees are not. The institutional reading is the best line the instrument owns: when the regulator confiscates the price, price competition does not die---it reappears as side payments, as it does in English bursaries under the fee cap, German Deutschlandstipendien, and the stipend wars of nominally free doctoral programmes. Means-blind admission is not money-blind students.

\paragraph{The post.} Beyond the shocks, the player---and only the player---receives dilemmas: roughly a dozen distinct cards over the library, offered with probability~0.3 a year from year two, each at most once per game. A card names a situation (a rival's professor is approachable; a donor offers money with conditions attached; an insurer offers a rider against the storm season), prints its price and its odds, and demands a decision before the year proceeds. Cards are drawn on the same separate random stream as the shocks, with all draws---including the success roll---consumed unconditionally, so no decision the player takes can perturb the applicant stream or the rivals' luck. An unanswered card resolves as a decline, which for the insurance-shaped cards means the gamble is taken by default: an unattended college buys no cover. Card prices are set against the investment technology $\gamma\sqrt{I}$, so that most cards are close to fair at a heavy investor's margins, poaching carries a zero-sum premium, and one card is priced deliberately below the technology. The reader who checks the arithmetic keeps their money.

\paragraph{The Founders' Reckoning.} The game's final table converts remaining endowments into league points: each college is credited $\min\!\big(\max(0, E)/15,\; 8\big)$ points, announced in the rules from year one and identical in both worlds. Two facts about the credit are worth separating. For the player it is dominated at any realistic scale: a final-round pound invested buys more quality than a pound reckoned until investment approaches five hundred per field, which no solvent college reaches. The Reckoning's bite lies elsewhere. The rival incumbents end typical games with hoards north of a thousand---capped, those hoards are worth up to eight points against a challenger, in both worlds. The Reckoning is a difficulty device wearing a savings product's clothes, and these notes say so plainly.

\paragraph{Events.} Each year each solvent college suffers an event with probability $\tfrac13$, drawn from a weighted table of thirty-one shocks: money shocks between $-26$ and $+32$, quality shocks between $-3$ and $+4$. Events are drawn from a random-number stream separate from the one generating applicants, so the two sources of luck never contaminate one another; money losses are truncated so that no event pushes an endowment below~5. Events fire before the year's league table is computed. The morning's news is news.

\paragraph{Endowments.} Starting endowments are world-specific, and deliberately so; they are also private, so the public starting position---the quality table---is identical across worlds. The market's bottom college holds a war chest (Table~\ref{tab:start}): the board recapitalises a turnaround, and the chest finances patient learning of the demand curve against the rent. Under the scheme the same chest would be poison to the calibration rather than to the college: with demand guaranteed by the clearing house, early over-spending becomes a winning strategy for the diligent and the asleep alike. The value of a war chest depends on the institution it is spent in.

\section{The first institution: the market}\label{sec:market}

In the market world a university chooses, for each field, a fee $\varphi_f \ge 0$ and an admission threshold $\tau_f$. Clearing is simultaneous and mechanical. Every applicant applies to every university whose fee for her field she can afford ($\varphi_f \le b$). Every university makes an offer to every applicant to a field whose score meets that field's bar ($s \ge \tau_f$). Every applicant holding at least one offer enrols at the offer maximising $T_f + \theta R_f$, ties broken at random; applicants without offers exit without consequence. Offers are unconditional: a university takes every student who accepts. Each enrolee beyond the eighth in a department costs $c = 24$---roughly three times a typical fee---so oversubscription is not a windfall but an accident, and occasionally an obituary. The fee cap sits at 25, above $c$, so overflow-for-profit is available in principle; a probe strategy built to try it (price just above $c$, bar on the floor, never ration) went bankrupt in over ninety per cent of seeds at every cost configuration tested and harvested essentially no overflow. To profit from a crowd one must first be crowded, and by the time a college is attractive enough to overflow at will, maintaining the quality that made it attractive has eaten the margin. The temptation exists; the customers do not.

Two features make this a game rather than an exercise. First, the demand curve is hidden: the player sees the cohort's score statistics but never its budgets, and must learn by experiment what the market will bear. The bottom college's central discovery is that it is a local monopolist over a captive segment---the students every rival rejects or prices out---and that the revenue-maximising fee for captives is interior: pricing at the floor buys almost no additional demand. Second, the accepted-offer rule creates a tail risk with a memorable signature. An offer book calibrated to a thirty per cent yield converts at ninety when a quality crossing redirects student choice, and the overage bill from one good year can end the college. The scripted sharp player forecasts its yield from the public league table and rations admission the moment it becomes attractive; human players learn the same lesson, usually the expensive way.

The stipend instrument of Section~\ref{sec:env} has a market-specific corollary worth recording: since it works only through affordability, it does nothing for a college whose fee is already within everyone's means. Merit aid in this market is an instrument of the expensive---a high list price for the captive segment, discounted at the gate for the bright---which is to say, second-degree price discrimination, and which is also to say, the American model. The bottom college, pricing low from necessity, has no use for it until it has a list price worth discounting.

One remark for the theorist. The market's clearing is itself a single round of student-proposing acceptance---students take their best offer---without quotas and without the right to say no. Seen this way, the two worlds differ in exactly two rights: the right to price, and the right to refuse.

\section{The second institution: the scheme}\label{sec:scheme}

In the scheme world the fee is fixed at $\bar{\varphi} = 5.5$ for every field and every college, admission is means-blind, and budgets are irrelevant. A university chooses, for each field, a quota $q_f \in \{0, 1, \ldots, 8\}$ and a threshold $\tau_f$. The seat costs of Section~\ref{sec:env} apply unchanged---the rent falls on all eight seats regardless---so the quota is a pure report to the clearing house: it costs nothing to state and binds only the match.

Matriculation is decided by student-proposing deferred acceptance (Gale and Shapley, 1962), run independently for each field. Students rank the four departments of their field strictly by $T_f + \theta R_f$, with vanishing random perturbations guaranteeing strictness. Departments hold responsive preferences: they rank students strictly by score, treat students below $\tau_f$ as unacceptable, and hold at most $q_f$. The algorithm is the standard one. Unmatched students propose down their lists; a department holds the best acceptable proposals up to its quota and rejects the rest, displacing previously held students where necessary; nothing is final until no rejected student has anywhere left to propose. The outcome is the student-optimal stable matching, unique under strict preferences and independent of the order in which proposals are processed (McVitie and Wilson, 1971). The implementation was audited mechanically: across three thousand randomly generated instances the produced matchings contain no blocking pair of either kind---no student preferring a department with a free acceptable seat, and no justified envy---and across twelve hundred simulated game rounds no quota is ever exceeded.

Two theoretical properties of this design do the game's work. The mechanism is strategy-proof for students (Roth, 1982)---fitting, since the students are honest automata---but no stable mechanism is strategy-proof for the colleges (Roth, 1985), and this one is manipulable in a specific, profitable direction: by capacities (S\"onmez, 1997). A department that reports fewer seats than its demand holds only its best proposers; its intake calibre rises, the $\kappa$ term feeds its teaching, and better teaching attracts better proposers. Since the rent is owed on the seats either way, the manipulation here is undiluted by economy: shrinking the report buys calibre at the price of forgone fees, exactly as the theory has it. The scripted sharp player nonetheless declines the move, and the reason is instructive: against rivals who spend their guaranteed income like rational actors (Appendix~\ref{app:ai}), the fees a small quota forgoes are worth more than the calibre it buys, and the calibrated best play is volume-first---every seat declared, the threshold doing the selecting. The manipulation is available, undiluted, and, at the margin, a mistake. Whether a manipulation pays is a statement about one's rivals as much as about the mechanism; a player who discovers both halves has discovered a theorem and its comparative statics.

The scheme also removes the market's tail risk. Deferred acceptance never over-fills a quota, so there is no overage and no yield lottery; what remains is the rent, which both worlds now share, and which the scheme's guaranteed demand makes easier to pay. Ruin under the scheme is reserved for the college that cannot fill what it keeps---and, as Section~\ref{sec:calib} shows, the same guarantee that spares the diligent also spares the asleep.

\section{Scripted play and calibration}\label{sec:calib}

Each world is calibrated with four scripted players run from the worst starting position over seeds $1$--$500$. The first three tiers are constructed from the same recipes in both worlds. \emph{Naive} makes static mid-range choices and spreads investment evenly. \emph{Sensible} tracks realised demand with its instruments, maintains its qualities against decay, and keeps a reserve. \emph{Sharp} exploits the institution: in the market, revenue hill-climbing on the captive segment, threshold-farming of the intake effect, yield forecasting and top-end rationing; in the scheme, volume-first admissions with the threshold as the selectivity instrument and high-bar stipends in the summit fight; in both, card purchases whenever liquidity and the printed odds allow---every card but the consultancy's. The fourth row, \emph{rustam}, is a benchmark contributed by the game's playtesting and run on every build: competent but unremarkable admissions, the full budget spent every year, split equally across the four qualities. Its interest lies in requiring no judgement at all---the quality technology is symmetric and concave, so equal division is close to optimal allocation, and what remains of the strategy is pure throttle. What pure throttle earns under each institution is a measurement worth publishing, and the harness additionally records whether it reaches the top by year five and holds it, which a well-calibrated scheme should not permit. Table~\ref{tab:calib} reports outcomes; Table~\ref{tab:diff} reports derived difficulty metrics on the same seeds.

\begin{table}[t]
\centering\small
\begin{tabular}{llccccc}
\toprule
 & & mean & \multicolumn{1}{c}{finishes} & improves & stagnates or & bankrupt \\
world & strategy & score & 1st (\%) & (\%) & declines (\%) & (\%) \\
\midrule
market & naive    & 0.24 & 0.4  & 23.6 & 76.4 & 99.6 \\
       & sensible & 0.37 & 2.0  & 30.0 & 70.0 & 10.0 \\
       & sharp    & 1.48 & 32.0 & 74.0 & 26.0 & 9.4 \\
       & rustam   & 0.81 & 12.4 & 50.6 & 49.4 & 61.0 \\
\addlinespace
scheme & naive    & 0.43 & 2.0  & 33.8 & 66.2 & 65.6 \\
       & sensible & 0.46 & 0.0  & 43.4 & 56.6 & 8.6 \\
       & sharp    & 2.17 & 33.4 & 91.8 & 8.2  & 1.2 \\
       & rustam   & 1.09 & 9.8  & 52.0 & 48.0 & 10.6 \\
\bottomrule
\end{tabular}
\caption{Calibration over 500 seeded histories per cell, player at the worst start (rank 4). Score $=$ initial rank $-$ final rank; ``improves'' is score $\ge 1$.}
\label{tab:calib}
\end{table}

\begin{table}[t]
\centering\small
\begin{tabular}{llcccc}
\toprule
world & strategy & s.d.\ of score & median escape year & never escapes (\%) \\
\midrule
market & naive    & 0.46 & 2    & 54 \\
       & sensible & 0.65 & 6    & 64 \\
       & sharp    & 1.19 & 3    & 9  \\
       & rustam   & 1.00 & 2    & 42 \\
\addlinespace
scheme & naive    & 0.68 & 10   & 66 \\
       & sensible & 0.54 & 11   & 56 \\
       & sharp    & 0.80 & 3    & 7  \\
       & rustam   & 1.11 & 8    & 46 \\
\bottomrule
\end{tabular}
\caption{Difficulty metrics, same 500 seeds. ``Escape'' is the first year the player leaves rank 4; medians are over escaping seeds.}
\label{tab:diff}
\end{table}

A caveat belongs up front, stated once. Comparing institutions through algorithm performance confounds how hard the world is with how good the hand-crafted algorithm happens to be in it. The naive and sensible rows are comparable by construction---the same abstract recipe mapped onto each world's instruments. The sharp rows are world-specific by design; they measure each institution's exploitability ceiling as probed with comparable effort, and nothing stronger is claimed.

With that caveat, four regularities stand. First, the rent bills the sleeping in both worlds, but only the market executes them: static play goes bankrupt in effectively every market seed (99.6 per cent), while under the scheme a third of static colleges survive the early grind---though they seldom prosper (2.0 per cent top the table), because rivals who spend their guaranteed income do not leave the summit to the idle. The institutional lesson sits where it belongs: where there is no price, there is no selection, and the clearing house keeps the asleep alive; it does not crown them. Second, mastery pays everywhere and calm is scheme-specific: sharp play reaches the top in a third of seeds in both worlds (32.0 and 33.4 per cent), but with much less noise under the scheme (s.d.\ 0.80 against 1.19, bankruptcy 1.2 against 9.4)---deferred acceptance abolishes the market's yield lottery, and what remains of sharp risk in the market is the accepted-offer rule meeting a quality crossing. Third, the rustam rows measure what pure throttle is worth under each institution: in the market it is a coin toss against ruin (bankrupt in 61 per cent of seeds), under the scheme it is safe and mediocre without being sovereign (tops 9.8 per cent, and reaches the summit by year five and holds it in none of five hundred)---full-throttle spending is a plan only where revenue cannot surprise you, and even there it is not a good one. Fourth, prudence buys survival, not distinction: sensible improves in 30 and 43 per cent of seeds but tops in 2.0 and 0.0---the Reckoning hands the hoarding incumbents up to eight points, and the reserve that keeps a sensible college solvent is precisely what it declines to spend on catching them.

Three probe strategies patrol the edges of the design, and their job is to fail. The overbooker of Section~\ref{sec:market} prices above $c$ and never rations: bankrupt in over ninety per cent of seeds, harvesting essentially no overflow. The scholar promises a stipend of 10 at a bar of 50 on an otherwise sensible game---nearly every matriculant qualifies---and is bankrupt in effectively every seed in both worlds: the stipend bill is a real route to ruin, and the instrument earns its keep only as designed, at a high bar in a summit fight, where it is decisive precisely when the fighting colleges sit within a stipend's width of one another in the students' eyes. The gambler accepts every card the post brings and does no better than sensible play once its extra bankruptcies are counted. An instrument that cannot be spammed is a calibrated one; these cannot.

\section{Discussion}\label{sec:disc}

The calibration is a designer's equilibrium, not nature's. The opponents are heuristics, tuned until the game's difficulty targets held; the numbers in Tables~\ref{tab:calib} and~\ref{tab:diff} characterise \game{}, not British higher education. The reader should know exactly what is, and is not, being claimed: the game does not adjudicate between the institutions, and its parameters were chosen for play, not estimated from anything.

What the comparison does teach is structural, and it survives the caveat. Remove the price instrument, and competition does not disappear---it migrates entirely into quality, where incumbents with hoarded endowments hold the high ground; the resulting reaction functions deter the half-hearted challenges that the market, indifferently, lets through, yet they cannot tell a compounding gambler from a compounding talent. Guarantee demand through a clearing house, and the fixed cost of seats stops selecting: the market forecloses on the passive, the scheme merely invoices them---and whether the invoiced ever reach the top depends not on the mechanism but on the conduct of their rivals. And the endowment experiment deserves its own line: the same war chest that finances patient learning under price competition finances indiscriminate expansion under guaranteed demand---which is why the two worlds start with different money and identical published tables. What capital is worth depends on the institution it is spent in.

For the classroom the game has one further use. Institutions are usually taught as rules and compared as theorems; here they are compared as lived experience, twenty minutes a world. A student who has died of overage understands the accepted-offer rule differently thereafter; a student who has paid three years' rent on empty seats holds views on the cost of ambition; a student who has tried the S\"onmez manipulation against rational rivals, and lost money doing it, has learned that a theorem about mechanisms is not yet a theorem about opponents; and a student who declines the consultancy's card after checking its odds against $\gamma\sqrt{I}$ has performed a cost--benefit analysis under a deadline, which is more than most syllabi manage. The theorems, of course, were always there. The game merely arranges for them to be discovered in anger.

\appendix

\section{Parameters and starting states}\label{app:params}

\begin{table}[h]
\centering\small
\begin{tabular}{llll}
\toprule
symbol & value & role & world \\
\midrule
$T$ & 20 & years per game & both \\
$n$ & 40 & applicants per year & both \\
--- & 8 & department capacity (cheap seats) & both \\
$\rho$ & 0.75 & annual rent per capacity seat, filled or not & both \\
$c$ & 24 & cost per enrolee beyond capacity & market \\
--- & 25 & fee cap (equals the budget truncation point) & market \\
$\bar\varphi$ & 5.5 & regulated fee & scheme \\
$r$ & 0.05 & interest on unspent funds & both \\
$\delta$ & 0.87 & quality retention & both \\
$\gamma$ & 1.5 & investment productivity, $\gamma\sqrt{I}$ & both \\
$\kappa$ & 0.14 & intake effect on teaching & both \\
$s$ & $N(60,13^2)$ on $[20,100]$ & school score & both \\
$b$ & $10 + 0.12(s{-}60) + N(0,3.8^2)$ on $[3,25]$ & budget & market \\
$\theta$ & $U[0,1.2]$ & research taste & both \\
$p_t$ & drift $0.02\tanh(\cdot/10)$ on $[0.3,0.7]$ & STEM share & both \\
--- & $\tfrac13$; weighted table & event probability; event law & both \\
--- & 5 & endowment floor under event losses & both \\
--- & 20; [40,95]; 1.0 & stipend cap; bar range (default 70); scheme pull per \pounds1k & both \\
--- & 15; 8 & Reckoning: endowment per league point; point cap & both \\
--- & 0.3; 2 & choice-card offer probability; first year offered & both \\
--- & 8; 1.35 & scheme AI cash float; scheme AI ceiling multiplier & scheme \\
\bottomrule
\end{tabular}
\caption{The full parameter set, as shipped in \texttt{DEFAULT\_PARAMS} inside the game file.}
\end{table}

\begin{table}[h]
\centering\small
\begin{tabular}{lccccccc}
\toprule
university & $E$ (market) & $E$ (scheme) & $R_S$ & $T_S$ & $R_H$ & $T_H$ & sum \\
\midrule
Harkness University    & 120 & 120 & 31 & 30 & 28 & 29 & 118 \\
Wexford Institute      & 100 & 100 & 36 & 20 & 28 & 17 & 101 \\
Millbrook Metropolitan & 130 & 95  & 18 & 24 & 16 & 25 & 83  \\
Greyfriars College     & 145 & 105 & 13 & 16 & 11 & 15 & 55  \\
\bottomrule
\end{tabular}
\caption{Starting states. Endowments are private and world-specific (Section~\ref{sec:env}); the published table---the qualities---is identical across worlds. The player is intended to take Greyfriars; the scoring quietly punishes a softer choice, since score is capped by starting rank minus one.}
\label{tab:start}
\end{table}

\section{The opponents}\label{app:ai}

The three rivals run fixed personalities under exactly the player's rules and information. \emph{Harkness}, the balanced incumbent, steers each quality toward a fixed target slightly above its starting level and otherwise hoards; in the scheme it defends the summit with a reaction function in quality: when the published gap to the college below narrows, its target tracks the challenger's total plus a margin---but the plan, defensive or not, is capped at 46 a year, so a challenger who can outspend the cap can outrun the defence. \emph{Wexford}, the prestige chaser, prices high (or, in the scheme, reports a small quota at a high bar), overweights research roughly two to one, and rescues its teaching only when collapse is visible; in the scheme it opens its war chest when displaced in the table. \emph{Millbrook}, the cash cow, prices low and fills every seat it can, invests a capped upkeep budget, and banks the rest; in the scheme it counter-invests from the hoard when overtaken or when a challenger closes to within four points. In the scheme, all three spend like the world they are in: with fee income guaranteed by the clearing house, precautionary saving is not caution but confusion, so their cash floats shrink to 8 (from the market's 15--20) and every spending ceiling above scales by 1.35---the balanced incumbent's defensive plan runs to 62 a year. Reserves calibrated against market risk exist only where market risk does. The scheme-only reaction functions exist because price competition does not: without them, any adaptive challenger strolled to the top; with them, and with the caps, the scheme blocks the middle of the skill distribution while remaining beatable from its top. All heuristics condition only on public information and their own books. Choice cards are the player's alone; the rivals' luck is confined to the shocks.

\section*{References}

\begin{description}
\item Gale, D., and L.\ S.\ Shapley (1962): ``College Admissions and the Stability of Marriage,'' \emph{American Mathematical Monthly}, 69(1), 9--15.
\item McVitie, D.\ G., and L.\ B.\ Wilson (1971): ``The Stable Marriage Problem,'' \emph{Communications of the ACM}, 14(7), 486--490.
\item Roth, A.\ E.\ (1982): ``The Economics of Matching: Stability and Incentives,'' \emph{Mathematics of Operations Research}, 7(4), 617--628.
\item Roth, A.\ E.\ (1985): ``The College Admissions Problem is not Equivalent to the Marriage Problem,'' \emph{Journal of Economic Theory}, 36(2), 277--288.
\item Roth, A.\ E., and M.\ Sotomayor (1990): \emph{Two-Sided Matching: A Study in Game-Theoretic Modeling and Analysis}. Cambridge University Press.
\item S\"onmez, T.\ (1997): ``Manipulation via Capacities in Two-Sided Matching Markets,'' \emph{Journal of Economic Theory}, 77(1), 197--204.
\end{description}

\end{document}
